% This must be in the first 5 lines to tell arXiv to use pdfLaTeX, which is strongly recommended.
\pdfoutput=1
% In particular, the hyperref package requires pdfLaTeX in order to break URLs across lines.
\documentclass[11pt]{article}
\usepackage{amsmath}
\usepackage{tabularx} % Add this in the preamble
\usepackage{amssymb}
% Remove the "review" option to generate the final version.
\usepackage{ACL2023}

% Standard package includes
\usepackage{times}
\usepackage{latexsym}

% For proper rendering and hyphenation of words containing Latin characters (including in bib files)
\usepackage[T1]{fontenc}
% For Vietnamese characters
% \usepackage[T5]{fontenc}
% See https://www.latex-project.org/help/documentation/encguide.pdf for other character sets

% This assumes your files are encoded as UTF8
\usepackage[utf8]{inputenc}

% This is not strictly necessary, and may be commented out.
% However, it will improve the layout of the manuscript,
% and will typically save some space.
\usepackage{microtype}

% This is also not strictly necessary, and may be commented out.
% However, it will improve the aesthetics of text in
% the typewriter font.
\usepackage{inconsolata}
\usepackage{graphicx} 
\usepackage{algorithm}
\usepackage{algorithmic}
\usepackage{float}
\usepackage{placeins}

% If the title and author information does not fit in the area allocated, uncomment the following
%
%\setlength\titlebox{<dim>}
%
% and set <dim> to something 5cm or larger.

\title{Seminar Report: Pixel-Level Certified Explanations via Randomized Smoothing}


\author{Daniya Niazi \\
\texttt{7062343,dani00003@stud.uni-saarland.de} \\
Universit{\"a}t des Saarlandes}

\date{}

\begin{document}
\maketitle
% \begin{abstract}
% 
% \end{abstract} 

\section{Introduction}

Deep neural networks achieve high accuracy in medical imaging, yet their decisions are often opaque. Post-hoc attribution methods produce heatmaps highlighting influential pixels, but these explanations can be unstable: small perturbations may drastically change the heatmap while the predicted class remains unchanged. This undermines trust and complicates comparison across methods. 

The paper \emph{Pixel-level Certified Explanations via Randomized Smoothing} proposes a framework that provides per-pixel robustness guarantees for any black-box attribution method \cite{anani2025pixel}. Instead of a single image-level score, it outputs certified pixel-wise labels: robustly important, robustly unimportant, or abstain.


\section{Background}
\subsection{Post-hoc Attribution Methods}
Attribution methods map an input image to a heatmap of pixel importance with respect to a predicted class. Three broad families dominate the literature: \textit{gradient-based} methods (e.g., saliency, Guided Backpropagation, Integrated Gradients); \textit{activation-based} methods (e.g., Grad-CAM variants); and \textit{perturbation-based} methods (e.g., Occlusion, RISE) . These methods are widely used because they are fast and easy to apply to trained models.

\subsection{Limitations of Post-hoc Attribution Methods}

Despite their popularity, post-hoc attribution methods exhibit fundamental robustness limitations. Prior work has shown that attribution maps can change substantially under small input perturbations, even when the classifier prediction remains unchanged.Ghorbani et al.~\cite{ghorbani2019fragile} demonstrated that adversarial perturbations can dramatically alter saliency maps while preserving the model output as shown in Figure \ref{fig:adversial_pertubation}, i.e., $f(x)=f(x')$ but $h(x)\neq h(x')$. Dombrowski et al.~\cite{dombrowski2019explanations} explained this instability by linking it to high curvature of decision boundaries, which causes sharp gradient fluctuations under small $\ell_2$ perturbations.
Subramanya et al.~\cite{subramanya2019fooling} showed that explanations can be intentionally manipulated to highlight arbitrary image regions without affecting the prediction, and Zhang et al.~\cite{zhang2020interpretable} demonstrated that models can be trained to produce misleading yet plausible explanations while maintaining predictive accuracy.
Kindermans et al.~\cite{kindermans2019reliability} further showed that certain attribution methods violate implementation invariance, changing under model reparameterizations that preserve the input-output function.

Kuppa and Le-Khac~\cite{kuppa2020blackbox} and Carmichael and Scheirer~\cite{carmichael2023unfooling} extended these concerns to model-agnostic explainers such as LIME and SHAP, demonstrating that perturbation-based explanations can also be destabilized or adversarially fooled.
\begin{figure}
    \centering
    \includegraphics[width=0.8\linewidth]{report_img/attr_pertub.png}
    \caption{Adversarial attack against feature-importance maps}
    \label{fig:adversial_pertubation}
\end{figure}

% Collectively, these works establish that explanation robustness does not follow from prediction robustness. Attribution maps can be fragile, manipulable, and sensitive to minimal perturbations, limiting their reliability in safety-critical domains.

\section{Existing Methods}

\subsection{Empirical Robustness Improvements}

Following the discovery that attribution maps are highly unstable, several works attempted to improve explanation robustness through empirical strategies. Chen et al.~\cite{chen2019robust} introduced attribution regularization during training, encouraging explanations to remain consistent under small perturbations of the input. By incorporating explanation stability directly into the loss function, their method reduced sensitivity of gradient-based attributions. Similarly, Dombrowski et al.~\cite{dombrowski2022towards} proposed training procedures that smooth gradients to mitigate geometric instability, and Wang et al.~\cite{wang2020smoothed} applied Gaussian smoothing to explanations to reduce high-frequency fluctuations.

These approaches improve stability in practice, but robustness remains empirical, method-dependent, and without formal guarantees on the perturbation radius under which explanations remain unchanged.

\subsection{Certified Attribution (Image-Level)}

To move beyond heuristic improvements, Levine et al.~\cite{levine2019certifiably} introduced the first certified robustness framework for explanations using randomized smoothing. Their method certifies a lower bound on the overlap between the top-$k$ most important features of clean and perturbed inputs within an $\ell_2$-bounded neighborhood. This provided the first probabilistic guarantee that important features remain stable under noise. Liu et al.~\cite{liu2022renyi} later strengthened these guarantees using Rényi differential privacy to obtain tighter bounds, and Wang and Kong~\cite{wang2023practical,wang2024certified} derived worst-case deviation bounds on attribution maps in norm space.

While these works provide formal certification, the guarantees remain image-level. Robustness is summarized by a single similarity score or norm bound per image, without identifying which individual pixels are provably stable. As a result, spatial interpretability of certified regions is still lacking.

\subsection{Segmentation Certification}

Fischer et al.~\cite{fischer2021scalable} extended randomized smoothing to structured outputs in segmentation tasks. Instead of certifying a single global prediction, their method certifies each output component independently, marking pixels as either robust or abstained if no guarantee can be provided. This demonstrates that pixel-wise certification is possible for high-dimensional outputs under $\ell_2$ perturbations.

However, attribution maps produce continuous importance scores rather than discrete class labels, making segmentation-style certification not directly applicable. Bridging this gap requires reformulating attribution as a segmentation problem to enable per-pixel robustness guarantees.


\section{Pixel-level Certified Explanations via Randomized Smoothing}

This section explains the workflow used to generate pixel-level certified attribution maps.
The method combines \textbf{sparsification} and \textbf{randomized smoothing} to determine
which pixels have attribution labels that are provably stable under bounded input perturbations.

\subsection{Problem Setting}
Let $f : \mathbb{R}^{C \times H \times W} \rightarrow \mathcal{Y}$ be a trained image classifier
and $h : \mathbb{R}^{C \times H \times W} \rightarrow \mathbb{R}^{N}$ be an attribution method
that produces a real-valued heatmap for an input image $x$, where $N = H \cdot W$.
Our goal is to certify, for each pixel independently, whether its attribution label
remains unchanged under small $\ell_2$ perturbations of the input.
\begin{figure}[h]
    \centering
    \includegraphics[width=1\linewidth]{report_img/certified_algo.png}
    \caption{Pixel-level certified explanations via random
ized smoothing.}
    \label{fig:placeholder}
\end{figure}


\subsection{Component I: Sparsification}
Attribution heatmaps are continuous‑valued and sensitive to noise. To obtain a discrete representation, the heatmap is sparsified. Given a heatmap $h(x)$ and a sparsification level $K$, the top $K\%$ of pixels with the highest attribution values are labeled as \emph{important} (1), and the remaining pixels are labeled as \emph{unimportant} (0). This produces a binary attribution map $h_K(x) \in \{0,1\}^N$. 

Let $h_f(x) = (a_1, \dots, a_N)$ be attribution scores and let $\mathrm{rank}(a_i)$ denote the descending rank of pixel $i$. For $K \in (0,100]$, define:

$$
h_K(x)_i =
\begin{cases}
1 & \text{if } \mathrm{rank}(a_i) \le \left\lceil \frac{K}{100} N \right\rceil, \\
0 & \text{otherwise}.
\end{cases}
$$
The parameter $K$ controls granularity: smaller values yield focused explanations, while larger values produce broader masks.
Thus, $h_K(x) \in \{0,1\}^N$ assigns each pixel to either the \emph{important} or \emph{unimportant} class.
Sparsification preserves the relative importance ranking of pixels while discarding unstable magnitude information. 

\subsection{Component II: Randomized Smoothing}

To assess robustness, the sparsified attribution is evaluated under noisy versions of the input image. Gaussian noise is added directly to the image,

$$
\epsilon \sim \mathcal{N}(0, \sigma^2 I),
$$

and the attribution pipeline (heatmap computation followed by sparsification) is recomputed for the perturbed input $x+\epsilon$. This process is repeated independently for $n$ samples.

\paragraph{Noisy attribution sampling.}
For a fixed pixel $i$, define the Bernoulli random variable

$$
Z_i = h_K(x+\epsilon)_i.
$$

\paragraph{Class probabilities.}
The probability that pixel $i$ belongs to class $c \in \{0,1\}$ under noise is

$$
p_c(i) = \mathbb{P}_{\epsilon \sim \mathcal{N}(0, \sigma^2 I)}\left(Z_i = c\right).
$$

\paragraph{Certified pixel label.}
Given a confidence threshold $\tau \in [0.5,1)$, the certified attribution map $\hat{h}_{\tau,K}(x) \in \{1,0,\varnothing\}^N$ is defined as

$$
\hat{h}_{\tau,K}(x)_i =
\begin{cases}
1 & \text{if } p_1(i) > \tau, \\
0 & \text{if } p_0(i) > \tau, \\
\varnothing & \text{otherwise}.
\end{cases}
$$

Pixels labeled $\varnothing$ are \emph{abstained} and receive no robustness guarantee.

\subsubsection{Interpretation of Abstention}
Abstention does not indicate that a pixel is unimportant.
It indicates that, under the chosen noise level and threshold,
there is insufficient statistical evidence to certify the pixel’s label.
This conservative behavior prevents misleading robustness claims.


\subsection{Certified Attribution Algorithm}


\begin{algorithm}[H]
\caption{Pixel-level Certified Attribution}
\label{alg:pixel_cert}
\begin{algorithmic}[1]
\REQUIRE Input image $x$, attribution method $h$, sparsification level $K$, noise level $\sigma$, threshold $\tau$, number of samples $n$
\ENSURE Certified attribution map $\hat{h}_{\tau,K}(x)$ and certified radius $R$

\STATE Initialize counters $c_1(i) \gets 0$ and $c_0(i) \gets 0$ for all pixels $i$
\FOR{$t = 1$ to $n$}
  \STATE Sample noise $\epsilon_t \sim \mathcal{N}(0,\sigma^2 I)$
  \STATE Compute heatmap $h(x+\epsilon_t)$
  \STATE Sparsify to binary mask $h_K(x+\epsilon_t)$
  \FOR{each pixel $i$}
    \IF{$h_K(x+\epsilon_t)_i = 1$}
      \STATE $c_1(i) \gets c_1(i) + 1$
    \ELSE
      \STATE $c_0(i) \gets c_0(i) + 1$
    \ENDIF
  \ENDFOR
\ENDFOR

\FOR{each pixel $i$}
  \STATE $p_1(i) \gets c_1(i)/n$, \quad $p_0(i) \gets c_0(i)/n$
  \IF{$p_1(i) > \tau$}
    \STATE $\hat{h}_{\tau,K}(x)_i \gets 1$
  \ELSIF{$p_0(i) > \tau$}
    \STATE $\hat{h}_{\tau,K}(x)_i \gets 0$
  \ELSE
    \STATE $\hat{h}_{\tau,K}(x)_i \gets \varnothing$
  \ENDIF
\ENDFOR

\STATE $R \gets \sigma \cdot \Phi^{-1}(\tau)$
\RETURN $\hat{h}_{\tau,K}(x), R$
\end{algorithmic}
\end{algorithm}


\subsection{Certified Robustness Guarantee}
The algorithm outputs a certified attribution map

$$
\hat{h}_{\tau,K}(x) \in \{1,0,\varnothing\}^N,
$$

where $\varnothing$ denotes abstention. For any non‑abstained pixel, the certified label is guaranteed invariant under any perturbation $\delta$ that satisfies

$$
\|\delta\|_2 \le R,\quad R = \sigma\,\Phi^{-1}(\tau).
$$

Here, $\|\delta\|_2$ is the Euclidean norm of the total input change, i.e., for a perturbed image $x'$, $\delta=x'-x$ and

$$
\|\delta\|_2 = \sqrt{\sum_j (x'_j - x_j)^2}.
$$

Thus, $R$ does not alter labels; it quantifies the **size of the robustness guarantee**: if the total input change is smaller than $R$, all certified 0/1 labels remain unchanged.

\subsection{Impact of Hyperparameters}
The method is governed by $(K,\sigma,\tau,n)$ and exhibits the following trade‑offs:
\begin{itemize}
\item \textbf{$K$ (sparsification)}: larger $K$ labels more pixels as important and tends to increase certifiability; smaller $K$ yields sharper focus but fewer certified pixels.
\item \textbf{$\sigma$ (noise)}: larger $\sigma$ increases $R$ (larger guarantee) but raises variability, often increasing abstention.
\item \textbf{$\tau$ (threshold)}: higher $\tau$ yields stricter certification and stronger guarantees, at the cost of reduced coverage.
\item \textbf{$n$ (samples)}: larger $n$ reduces estimation variance and supports stricter $\tau$, but increases computation.
\end{itemize}
The framework is attribution‑agnostic and applies to any black‑box attribution method that produces pixel‑wise scores.



\subsection{Evaluation Framework}
The authors evaluation across three metrics. Each one answers a different question: (i) how many pixels are certifiably stable, (ii) whether certified important pixels fall in the correct region, and (iii) whether removing certified important pixels hurts the prediction. All metrics are computed from the certified map $\hat{h}_{\tau,K}(x)$ and the model output.

\subsubsection{\% Certified (coverage)}
This metric measures the amount of the image that is certified (i.e., not abstained). 
Let $N=H\cdot W$ and use $\varnothing$ for abstention. Then

$$
\%\mathrm{certified}(x) = \frac{|\{i \mid \hat{h}_{\tau,K,i}(x) \neq \varnothing\}|}{N}.
$$
For each image, count all pixels labeled 0 or 1, divide by $N$, then average over images. A higher value means the method certifies more pixels under the chosen $(\sigma,\tau,n)$.

\subsubsection{Certified Grid Pointing Game (Certified GridPG)}
This metric checks localization on grid datasets (images split into $m\times m$ cells), based on GridPG \cite{bohle2021convolutional}. Suppose the correct cell is $c$ (In this cases the truth cell is fixed i,e cell 0). Let $x_i$ be the pixels in cell $i$, and define $A_c(p)=1$ if pixel $p$ is certified important (label 1), otherwise $0$. The certified localization score is

$$
cL_c = \frac{\sum_{p \in x_c} A_c(p)}{\sum_{j=1}^{m^2} \sum_{p \in x_j} A_c(p)}.
$$
\begin{figure}
    \centering
    \includegraphics[width=1\linewidth]{report_img/grid.png}
    \caption{Certified Grid Example}
    \label{fig:placeholder}
\end{figure}

For each image, sum certified-1 pixels inside the correct cell and divide by all certified-1 pixels in the image. If there are no certified-1 pixels, the image is skipped (or assigned 0 by convention). The final score is the mean across images.

\subsubsection{Faithfulness via Certified Deletion}
Faithfulness asks if the certified important pixels are truly used by the model. Let $S$ be the set of certified-1 pixels. Create a perturbed image $x^{\setminus S}$ by deleting or masking pixels in $S$ (e.g., set them to 0). The faithfulness score is the confidence drop for the predicted class $y$:

$$
\Delta f(x) = f_y(x) - f_y\big(x^{\setminus S}\big).
$$
For each image, remove the certified-1 pixels and measure how much the model’s confidence decreases. Larger drops indicate that the certified pixels carry predictive evidence.

The paper reports these three metrics across multiple attribution methods to show the trade-off between coverage (robustness), localization, and causal relevance.

% \begin{table}[h!]
% \centering
% \caption{Summary of the three evaluation metrics used for certified attributions.}
% \begin{tabularx}{\linewidth}{|X|X|X|}
% \hline
% \textbf{Metric} & \textbf{What it measures} & \textbf{Why it matters} \\
% \hline
% \% certified & Fraction of pixels certified as 0 or 1 & Direct robustness indicator; higher is more stable \\
% \hline
% Certified GridPG & Certified localization on grid images & Tests whether robust pixels align with the correct subimage \\
% \hline
% Faithfulness (deletion) & Confidence drop after removing certified top-$K$ pixels & Evaluates whether certified pixels truly support the prediction \\
% \hline
% \end{tabularx}
% \label{tab:probmodel}
% \end{table}


% \section{Visualization Placeholder}
% \begin{figure}[h]
%     \centering
%     \includegraphics[width=1\linewidth]{.png}
%     \caption{Placeholder for a certified attribution visualization (e.g., overlayed top-$K$ certified pixels).}
%     \label{fig:enter-label}
% \end{figure}

\section{Implementation and Experiments}

\subsection{Experimental Setup}
ImageNet; ResNet-18, Wide ResNet-50-2, ResNet-152, VGG-11, ViT-B/16; $224\times224$ images; 100 high-confidence samples. Twelve attribution methods across backprop/activation/perturbation families (ViT omits LRP, CAM). Input and final spatial layers; deeper maps are upsampled. Defaults: $\sigma=0.15$, $n=100$, $\tau=0.75$, $\alpha=0.001$, so $R\approx0.10$. GridPG uses 2$\times$2 grids (448$\times$448), 100 grids per method.

\subsection{Result Analysis}
LRP and RISE are the most consistent across robustness, localization, and faithfulness. In Figure \ref{fig:paper_fig5}, they retain relatively high \% certified, especially at the input layer.

\begin{figure}[h]
    \centering
    \includegraphics[width=1\linewidth]{report_img/per_certified_fig_5.png}
    \caption{Comparison of the per-pixel certification rate }
    \label{fig:paper_fig5}
\end{figure}

In Figure \ref{fig:paper_fig6}, they also show strong Certified GridPG. In Figure~8, they cause steep confidence drops under deletion, indicating higher faithfulness. Together, these figures show they are both stable and meaningful.

\begin{figure}[h]
    \centering
    \includegraphics[width=1\linewidth]{report_img/per_grid_fig_6.png}
    \caption{Comparison of the certified localization}
    \label{fig:paper_fig6}
\end{figure}



The robustness--localization trade-off appears when comparing Figures \ref{fig:paper_fig5} and \ref{fig:paper_fig6}. Some methods are highly certified but poorly localized, so robustness alone is insufficient. Figure \ref{fig:paper_fig7} summarizes this balance.
\begin{figure}[h]
    \centering
    \includegraphics[width=1\linewidth]{report_img/robvsloc.png}
    \caption{robustness-localization trade-off}
    \label{fig:paper_fig7}
\end{figure}


\subsubsection{Relationship Between $K$ and Layer Depth on \% Certified}
Figures \ref{fig:paper_fig2} and \ref{fig:paper_fig5} show that smaller $K$ yields fewer, more focused certified pixels, while larger $K$ yields broader regions.
\begin{figure}[h]
    \centering
    \includegraphics[width=0.8\linewidth]{report_img/sparcification.png}
    \caption{Certified Pixel at Different Sparsification Level (k)}
    \label{fig:paper_fig2}
\end{figure}

At the input layer, small $K$ is unstable under noise and increases abstention; larger $K$ stabilizes the top set. At the final layer, smaller $K$ isolates discriminative regions and improves localization. Overall: larger $K$ helps input-layer stability; smaller $K$ helps deep-layer sharpness.

\subsection{Outcome}
Qualitative maps show smaller $K$ gives finer details, larger $K$ gives broader regions. Figure \ref{fig:paper_fig5} shows \% certified decreases as $R$ increases. Figure \ref{fig:paper_fig6} shows Certified GridPG improves when certified pixels concentrate in the correct cell. Faithfullness evaluation also shows that deletion sharply reduces confidence for LRP/RISE, matching their robustness and localization.



\section{Limitations}
Sparsification discards magnitude information, so graded relevance can be lost. The method is also computationally heavy because certification requires many noisy samples per image. Its guarantees apply only to $\ell_2$-bounded noise and may not cover broader distribution shifts. Finally, stricter thresholds increase abstention, reducing certified coverage.

\section{Reflection}
Recasting attribution as segmentation is elegant and yields pixel-wise guarantees. The evaluation is strong because it separates robustness, localization, and faithfulness. My main concern is sampling: $n=100$ may be fine for $224\times224$ but could be insufficient for higher-resolution or fine-detail domains, while increasing $n$ is costly. Some trends are expected (e.g., smaller $K$ means fewer certified pixels); the more informative findings are the layer-wise and robustness--localization trade-offs. Overall, the framework is promising but faces scalability and sparsification limits.

% \section*{Acknowledgements}
% This report is based on the research paper \textit{Pixel-level Certified Explanations via Randomized Smoothing}. I thank the authors for their comprehensive work and Professor Jonas for supervision, which served as the foundation of this seminar report.

\bibliography{anthology,custom}
\bibliographystyle{acl_natbib}

\end{document}
